%%=============================================================================
%% Inleiding
%%=============================================================================

\chapter{\IfLanguageName{dutch}{Inleiding}{Introduction}}%
\label{ch:inleiding}

Goossens Warmte- en Koeltechniek is een Belgisch HVAC-bedrijf dat Installware gebruikt ter ondersteuning van de buitendienst.
Deze bachelorproef richt zich op de registratie en verwerking van digitale werkbongegevens binnen deze bedrijfscontext.
De doelgroep voor deze bachelorproef zijn de buitendienstmedewerkers en de kantoormedewerkers die deze gegevens
raadplegen en verwerken.

\section{\IfLanguageName{dutch}{Probleemstelling}{Problem Statement}}%
\label{sec:probleemstelling}

Wanneer medewerkers gegevens aanpassen zonder betrouwbare netwerkverbinding,
moet duidelijk zijn hoe deze gegevens lokaal worden bewaard en hoe ze worden
verwerkt zodra er opnieuw een netwerkverbinding beschikbaar is.

Het onderzoek vergelijkt twee concreet uitgewerkte aanpakken:
Een lokale wachtrij met centrale conflictcontrole en een aanpak met conflictvrije gerepliceerde datatypes.
Het wordt afgebakend tot werknotities en materiaalregistraties. Tijdens de analyse van de huidige werkwijze
wordt getoetst of deze gegevens en de bijbehorende gebruiksscenario’s representatief zijn voor het bedrijf.
De technische vergelijking die gemaakt zal worden bij de proof of concept zal geen toegang tot de broncode van Installware nodig hebben.

\section{\IfLanguageName{dutch}{Onderzoeksvraag}{Research question}}%
\label{sec:onderzoeksvraag}

De hoofdonderzoeksvraag: “Welke van twee synchronisatieaanpakken, een lokale wachtrij met centrale conflictcontrole of een samenvoeging op basis van conflictvrije gerepliceerde datatypes, biedt de beste afweging tussen gegevenscorrectheid, synchronisatietijd en implementatiecomplexiteit voor offline bewerkte werknotities en materiaalregistraties op digitale werkbonnen bij Goossens Warmte- en Koeltechniek?”

Deelvragen probleemdomein:
(1) Hoe registreren en verwerken buitendienstmedewerkers en kantoormedewerkers momenteel werknotities en materiaalregistraties met Installware?
(2) Welke mogelijkheden en beperkingen ervaren zij bij het raadplegen, bewerken en synchroniseren van deze gegevens tijdens en na een onderbroken netwerkverbinding?
(3) In welke situaties wijzigen meerdere medewerkers dezelfde werkbon en bestaan daar al bepaalde regels voor?
(4) Welke functionele en technische vereisten stelt het bedrijf aan offline gegevensvastlegging en de daaropvolgende synchronisatie?

Deelvragen oplossingsdomein:
(1) Hoe kunnen beide synchronisatieaanpakken worden ingericht voor de geselecteerde werkbongegevens en bedrijfsvereisten?
(2) Hoe verwerken beide implementaties wijzigingen bij verbindingsuitval, onderbroken synchronisatie, herhaalde verzending en gelijktijdige bewerking?
(3) Hoe verschillen beide implementaties in gegevenscorrectheid en synchronisatietijd bij dezelfde gegevens, gebruiksscenario’s en netwerkomstandigheden?
(4) Welke implementatiecomplexiteit en beperkingen brengt elke aanpak met zich mee?

\section{\IfLanguageName{dutch}{Onderzoeksdoelstelling}{Research objective}}%
\label{sec:onderzoeksdoelstelling}

Het doel van deze bachelorproef is een onderbouwd advies te formuleren over de geschiktheid van beide synchronisatieaanpakken
voor de bedrijfsvereisten. Het onderzoek levert een analyse van de huidige werkwijze en vereisten, een zelfstandig prototype met
twee vergelijkbare implementaties en een evaluatie op basis van reproduceerbare tests op.
Gegevenscorrectheid wordt daarbij beoordeeld aan de hand van zowel het behoud van registraties als de naleving
van de vastgestelde regels voor werkbonnen. Het onderzoek heeft niet als doel een volledige vervanging van Installware,
maar een technisch onderbouwde vergelijking die het bedrijf inzicht geeft in de mogelijkheden en beperkingen van beide aanpakken.

\section{\IfLanguageName{dutch}{Opzet van deze bachelorproef}{Structure of this bachelor thesis}}%
\label{sec:opzet-bachelorproef}

% Het is gebruikelijk aan het einde van de inleiding een overzicht te
% geven van de opbouw van de rest van de tekst. Deze sectie bevat al een aanzet
% die je kan aanvullen/aanpassen in functie van je eigen tekst.

De rest van deze bachelorproef is als volgt opgebouwd:

In Hoofdstuk~\ref{ch:stand-van-zaken} wordt een overzicht gegeven van de stand van zaken binnen het onderzoeksdomein, op basis van een literatuurstudie.

In Hoofdstuk~\ref{ch:methodologie} wordt de methodologie toegelicht en worden de gebruikte onderzoekstechnieken besproken om een antwoord te kunnen formuleren op de onderzoeksvragen.

% TODO: Vul hier aan voor je eigen hoofstukken, één of twee zinnen per hoofdstuk

In Hoofdstuk~\ref{ch:conclusie}, tenslotte, wordt de conclusie gegeven en een antwoord geformuleerd op de onderzoeksvragen. Daarbij wordt ook een aanzet gegeven voor toekomstig onderzoek binnen dit domein.